\scanpage{68}
\mspage{74}

On vérifie
\[
\left\{\begin{aligned}
\unsure{\mathscr D|\mathscr P}(X)_{/x}&\simeq\mathscr D(X_{/x}) & x&\in\mathrm{Ob}\,X\\
\mathscr D(X)_{/x^\circ}&\simeq\mathscr D({}_{x\backslash}X) & x^\circ&\in\mathrm{Ob}\,X^\circ
\end{aligned}\right.
\]

Je dis que $p_X$, \inserted{$q_X$} \uline{sont des $\mathbb D$-équiv. strictes}, \inserted{et costrictes} \deleted{\unsure{$q_X$ une} $\mathbb D$-équivalence costricte}. Par les formules précédentes, \deleted{\unsure{on est ramené}} cela signifie que les $\mathscr D(X_{/x})$, $\mathscr D({}_{x\backslash}X)$ sont $\deleted{\unsure{d}}\mathbb D$-contractiles. Donc que si $Z$ a un objet final, $\mathscr D(Z)$ est $\mathbb D$-contractile.

Mais considérons
\diagram{\begin{tikzcd}[ampersand replacement=\&]
\mathscr D(Z) \arrow[dr,"q_Z"] \& \\
\& Z^\circ,
\end{tikzcd}}
il suffit \deleted{\unsure{que}} \unsure{de} prouver que c'est une \inserted{$\mathbb D$-}équiv. stricte (puisque $Z^\circ$, ayant un objet initial, est $\mathbb D$-contractile), donc que les $ {}_{z_{\unsure{0|\circ}}\backslash}\mathscr D(Z)$ sont $\mathbb D$-contractiles

On a par la formule précédente
\[
{}_{z_0\backslash}\mathscr D(Z)\simeq\mathscr D({}_{z_0\backslash}Z)=\mathscr D(T),
\]
où cette fois $T$ a un objet initial \inserted{$\unsure{o|0}$} et \deleted{final.} un objet final $e$. Considérons alors l'objet
\diagram{\begin{tikzcd}[ampersand replacement=\&]o\arrow[r]\& e\end{tikzcd}}
de $\mathscr D(T)$, c'est un objet final, car \deleted{\unsure{$\forall$}} pour tout $x\longrightarrow y$ de $\mathscr D(T)$, il y a une unique flèche
\diagram{\begin{tikzcd}[ampersand replacement=\&]
x\arrow[r] \& y\arrow[d]\\
o\arrow[u]\arrow[r] \& e
\end{tikzcd}}
% archive sheet number: 67
